\subsection{群 SL(2, k) 的线性表示}

回顾（《代数学》第三章 §8 第 9 节）我们用 \(\mathbf{SL}(2,k)\) 表示行列式等于 1 的、系数在 \(k\) 中的 2 阶方阵组成的群。若 \(x\in \mathfrak{sl}(2,k)\) 幂零，则 \(x^{2} = 0\) （《代数学》第七章 §5 命题 5 的推论 3），且 \(e^x = 1 + x\in \mathbf{SL}(2,k)\) 。若 E 是有限维向量空间而 \(\rho\) 是 \(\mathfrak{sl}(2,k)\) 在 E 上的线性表示，则 \(\rho (x)\) 幂零，于是 \(e^{\rho (x)}\) 有定义（第一章 §6 第 3 节）。

定义 2. 设 E 是有限维向量空间，\(\rho\) （相应地 \(\pi\) ）是 \(\mathfrak{sl}(2,k)\) （相应地 \(\mathbf{SL}(2,k)\) ）在 E 上的线性表示。称 \(\rho\) 与 \(\pi\) 相容，若对 \(\mathfrak{sl}(2,k)\) 的每个幂零元素 x，\(\pi(e^{x}) = e^{\rho(x)}\) 。

换言之，\(\rho\) 与 \(\pi\) 相容，若对 \(\mathfrak{sl}(2,k)\) 的每个幂零元素 x，\(\rho\) 在 kx 上的限制与 \(\pi\) 在群 \(1 + kx\) 上的限制相容（第七章 §3 第 1 节）。

若 \(\rho\) 与 \(\pi\) 相容，则 \(\rho\) 与 \(\pi\) 的对偶表示、m 次张量幂以及 m 次对称幂也分别相容（第七章 §5 第 4 节 引理 1 (i) 与 (ii)）。对 \(\rho\) 与 \(\pi\) 在 \(\rho\) 与 \(\pi\) 下稳定的向量子空间上诱导的表示，同样成立（同前引）。

特别地，表示 \(\rho_{m}\) （\(\mathfrak{sl}(2,k)\) 在 \(V(m)\) 上的，见第 3 节）与 \(\pi_{m}\) （恒同表示 \(\pi_{1}\) （\(\mathbf{SL}(2,k)\) 的）的 m 次对称幂）相容。如上置 \(e_{n}^{(m)} = \binom{m}{n} u^{m-n} v^{n}\) ，我们有

\[
\pi_ {m} (s) e _ {n} ^ {(m)} = \binom {m} {n} (s u) ^ {m - n} (s v) ^ {n}\tag{3}
\]

对 \(s \in \mathbf{SL}(2, k)\) 与 \(0 \leq n \leq m\) 成立。

定理 2. 设 \(\rho\) 是 \(\mathfrak{sl}(2,k)\) 在有限维向量空间 E 上的线性表示。

\begin{enumerate}
\def\labelenumi{(\roman{enumi})}
\item
  存在唯一的线性表示 \(\pi\) （\(\mathbf{SL}(2,k)\) 在 E 上的），它与 \(\rho\) 相容。
\item
  向量子空间 \(\mathrm{F}\) （含于 \(\mathrm{E}\) 的）在 \(\pi\) 下稳定，当且仅当它在 \(\rho\) 下稳定。
\item
  设 \(x \in \mathrm{E}\) 。则 \(\pi(s)x = x\) 对一切 \(s \in \mathbf{SL}(2,k)\) 成立，当且仅当 \(x\) 在 \(\rho\) 下不变（即 \(\rho(a)x = 0\) 对一切 \(a \in \mathfrak{sl}(2,k)\) 成立）。
\end{enumerate}

\(\pi\) 的存在性由前述及定理 1 得出。另一方面，我们知道群 \(\mathbf{SL}(2,k)\) 由形如以下的元素生成：

\[
e ^ {t X _ {+}} = \left( \begin{array}{c c} 1 & t \\ 0 & 1 \end{array} \right) \quad e ^ {- t X _ {-}} = \left( \begin{array}{c c} 1 & 0 \\ t & 1 \end{array} \right)
\]

其中 \(t \in k\) （《代数学》第三章 §8 第 9 节 命题 17）。这就证明了 \(\pi\) 的唯一性。

断言 (ii) 与 (iii) 由前述以及第七章 §3 第 1 节 引理 1 (i) 得出。证毕。

因此每个有限维 \(\mathfrak{sl}(2,k)\) -模都有唯一的 \(\mathbf{SL}(2,k)\) -模结构，称之为与其 \(\mathfrak{sl}(2,k)\) -模结构相伴的结构。

注. 当 k 是 R、C 或完备非离散超度量域时，\(\mathfrak{sl}(2,k)\) 是 \(\mathbf{SL}(2,k)\) 的李代数。设 \(\rho\) 与 \(\pi\) 如定理 2。同态 \(\pi\) 是从 \(\mathbf{SL}(2,k)\) 到 \(\mathbf{GL}(\mathrm{E})\) 的李群同态：当 \(\mathrm{E} = \mathrm{V}(m)\) 时这是显然的，一般情形由定理 1 得出。由第七章 §3 第 1 节，\(\rho(X_{+}) = \mathrm{L}(\pi)(X_{+})\) ，\(\rho(X_{-}) = \mathrm{L}(\pi)(X_{-})\) 。故 \(\rho = \mathrm{L}(\pi)\) （其逆见习题 18）。

命题 5. 设 E、F 是有限维 \(\mathfrak{sl}(2,k)\) -模，并设 \(f \in \operatorname{Hom}_{k}(\operatorname{E},\operatorname{F})\) 。下列条件等价：

\begin{enumerate}
\def\labelenumi{(\roman{enumi})}
\item
  \(f\) 是 \(\mathfrak{sl}(2,k)\) -模的同态；
\item
  \(f\) 是 \(\mathbf{SL}(2,k)\) -模的同态。
\end{enumerate}

条件 (i) 意味着 \(f\) 是 \(\mathfrak{sl}(2,k)\) -模 \(\mathrm{Hom}_k(\mathrm{E},\mathrm{F})\) 的不变元，条件 (ii) 意味着 \(f\) 是 \(\mathbf{SL}(2,k)\) -模 \(\mathrm{Hom}_k(\mathrm{E},\mathrm{F})\) 的不变元。由于这两种模结构依第七章 §5 第 4 节 引理 1 (iii) 是相伴的，本命题由定理 2 (iii) 得出。

定义 3. 群 \(\mathbf{SL}(2,k)\) 的伴随表示是 \(\mathbf{SL}(2,k)\) 在 \(\mathfrak{sl}(2,k)\) 上由下式定义的线性表示 Ad：

\[
\operatorname{Ad} (s). a = s a s ^ {- 1}
\]

对一切 \(a \in \mathfrak{sl}(2, k)\) 与一切 \(s \in \mathbf{SL}(2, k)\) 成立。

当 k 是 R、C 或完备非离散超度量域时，我们重新得到第三章 §3 第 12 节 定义 7（参见该处命题 49）。

由第七章 §5 第 4 节 引理 1 (i) 与 (ii)，\(\mathfrak{sl}(2,k)\) 与 \(\mathbf{SL}(2,k)\) 的伴随表示相容。由第七章 §3 第 1 节 注 2，\(\operatorname{Ad}(\mathbf{SL}(2,k)) = \operatorname{Aut}_e(\mathfrak{sl}(2,k))\) 。

